% Part: first-order-logic
% Chapter: syntax-and-semantics
% Section: introduction

\documentclass[../../../include/open-logic-section]{subfiles}

\begin{document}

\olfileid{fol}{syn}{int}

\olsection{جان پچھان}

پہلے درجے دی منطق دا نظریہ تے اوہدے بارے اُتلا نظریہ بناؤن لئی سانوں پہلاں ایہدیاں عبارتاں دی بناوٹ تے معنی دی تعریف کرنی پئے گی۔ پہلے درجے دی منطق دیاں عبارتاں ٹرماں تے !!{formula}s نیں۔ ٹرم !!{variable}s، !!{constant}s تے !!{function}s توں بندے نیں۔ اگّوں !!^{formula}s، !!{predicate}s نوں ٹرماں نال جوڑ کے بندے نیں (ایہ سبھ توں نِکّے، «جوہری» !!{formula}s نیں)، تے فیر جوہری !!{formula}s توں منطقی جوڑاں تے مسوّراں دی مدد نال ہور گنجل دار فارمولا بنا سکدے آں۔ بناوٹ دے قاعدے دَسن دے کئی وکھ وکھ ڈھنگ نیں۔ اسیں بس اک ممکن ڈھنگ دِندے آں۔ دوجے نظام وکھ وکھ نشان چُنن گے، بنیادی جوڑاں دے وکھ وکھ سیٹ رکھن گے، قوساں وکھرے طریقے نال ورتن گے (یا بالکل نہیں ورتن گے، جیویں پولش لکھت وچ)۔ پر سارے طریقیاں دی سانجھی گل ایہ اے کہ بناوٹ دے قاعدے ٹرماں تے !!{formula}s دے سیٹ دی \emph{استقرائی} تعریف کردے نیں۔ جے ایہ ٹھیک طرح ہووے، تاں بناوٹ دے قاعدیاں مطابق ہر عبارت بنیادی طور اُتے بس اک ای ڈھنگ نال بن سکدی اے۔ عبارت دی ایہ استقرائی تعریف، جیہدے نال اوہ \emph{اکّو ڈھنگ نال پڑھی جا سکدی} اے، سانوں ایہناں عبارتاں دے معنی وی اوسے ڈھنگ نال، یعنی استقرائی تعریف نال، دین دی کھُل دیندی اے۔

عبارتاں نوں معنی دینا علمِ معنی دا کم اے۔ علمِ معنی دا مرکزی خیال !!a{structure} وچ پورا اترنا اے۔ !!^a{structure} زبان دیاں بنیادی اکائیاں نوں معنی دیندی اے۔ !!a{domain} شےواں دا غیر خالی سیٹ اے۔ مسوّراں دی تعبیر ایویں ہوندی اے کہ اوہ ایس دائرے دیاں شےواں اُتے پھردے نیں۔ !!{constant}s نوں دائرے دے رکن دِتّے جاندے نیں، !!{function}s نوں !!{domain} توں اوسے دائرے ول فنکشن دِتّے جاندے نیں، تے !!{predicate}s نوں !!{domain} اُتے تعلق دِتّے جاندے نیں۔ !!{domain} تے بنیادی لفظاں لئی دِتّیاں تعبیراں رَل کے !!a{structure} بناؤندیاں نیں۔ !!^{variable}s، !!{formula}s وچ آ سکدے نیں، تے معنی دین لئی سانوں !!{element}s وی !!{domain} توں ایہناں متغیراں نوں دین پَین گے۔ ایہ متغیر قدر بندی اے۔ آخر پورا اترن دا تعلق ایہناں گلّاں نوں جوڑدا اے۔ !!^a{formula}، !!a{structure}~$\Struct{M}$ وچ متغیر قدر بندی~$s$ دے لحاظ نال پورا اُتر سکدا اے، جیہنوں $\Sat{M}{!A}[s]$ لکھدے آں۔ ایس تعلق دی وی~$!A$ دی بناوٹ اُتے استقرا نال تعریف ہوندی اے۔ منطقی جوڑاں دیاں سچائی سارنیاں ورت کے، مثال لئی، $!A \land !B$ دا پورا اترنا، $!A$ تے ~$!B$ دے پورا اترن (یا نہ اترن) نال تعریف ہوندا اے۔ فیر پتا لگدا اے کہ جے !!{formula}~$!A$ اک !!a{sentence} اے، یعنی اوہدے وچ کوئی آزاد متغیر نہیں، تاں متغیر قدر بندی دا کوئی اثر نہیں۔ سو اسیں !!{sentence}s دے !!{structure}s وچ بس پورا اترن (یا نہ اترن) دی گل کر سکدے آں۔

جملیاں لئی پورا اترن دے تعلق $\Sat{M}{!A}$ دی بنیاد اُتے اسیں معنی دیاں بنیادی اصطلاحاں، صحت، لازمی نتیجہ تے پورا اتر سکنا، دی تعریف کر سکدے آں۔ اک جملہ صحیح اے، $\Entails !A$، جے ہر ساخت وچ پورا اُترے۔ اوہ !!{sentence}s دے کسے سیٹ دا لازمی نتیجہ اے، $\Gamma \Entails !A$، جے ہر !!{structure} جیہدے وچ سارے !!{sentence}s، جیہڑے~$\Gamma$ وچ نیں، پورے اُترن، اوہدے وچ~$!A$ وی پورا اُترے۔ جملیاں دا سیٹ پورا اُتر سکدا اے جے کوئی !!{structure} ایہو جیہی ہووے جیہدے وچ اوہدے سارے !!{sentence}s اکٹھے پورے اُترن۔ کیوں جے !!{formula}s دی استقرائی تعریف ہوندی اے تے اگّوں پورا اترن دی تعریف !!{formula}s دی بناوٹ اُتے استقرا نال ہوندی اے، اسیں استقرا ورت کے اپنے علمِ معنی دیاں خاصیتاں ثابت کر سکدے آں تے ایہناں معنی دیاں اصطلاحاں دا آپس دا تعلق وکھا سکدے آں۔

\end{document}
